\documentclass[10pt,a4paper]{article}
\usepackage[margin=0.8in]{geometry}
\usepackage{times}
\usepackage{enumitem}
\usepackage[colorlinks,allcolors=blue]{hyperref}
\setlist{leftmargin=1.3em,topsep=1pt,itemsep=1pt,parsep=0pt}
\setlength{\parindent}{0pt}
\setlength{\parskip}{2.5pt}

\title{\vspace{-1.5em}AI Usage Appendix\\[2pt]
\large LSTM vs.\ Transformer Encoder for Text Classification (AG News)}
\author{Jihun Jang (ETRI School, 02521122) \quad Deep Learning 2, 2026}
\date{}

\begin{document}
\maketitle
\vspace{-1.5em}

\section*{1.\ AI tools used}
Claude (Anthropic, in the Claude Code environment) was used as a coding assistant and learning tutor.

\section*{2.\ Where AI was used}
\begin{itemize}
\item Concept explanations: LSTM recurrence and gates, self-attention and multi-head attention, positional encoding, masking, overfitting, data efficiency, and so on.
\item Code drafts: the preprocessing pipeline, both models, the training and evaluation loops, the ablation, gradient saliency and attention analysis, and figure generation.
\item Debugging, support in interpreting results, and drafting the plan and report.
\item The process emphasized internalizing concepts through a tutor workflow --- predict, run, get feedback, and recall at each decision point --- rather than simple ghost-writing.
\end{itemize}

\section*{3.\ What the team corrected or supplemented}
\begin{itemize}
\item Dropout was controlled to start from a common 0.1 (instead of arbitrary asymmetric values of 0.3 and 0.1) and then tuned by validation.
\item The extension topic was changed from the AI's initial active-learning link to a Transformer mechanism analysis.
\item The writing style of the plan and report was tightened for concision and accuracy, removing unnecessary phrasing.
\end{itemize}

\section*{4.\ Decisions made by the team}
\begin{itemize}
\item The choice of the ablation topic (data efficiency), and the demand for a more insightful alternative.
\item The scope and direction of the extension, and the deliverable structure and writing style.
\item Hypotheses and predictions before each experiment (e.g., which model drops more when data is reduced), and the interpretation of results.
\end{itemize}

\section*{5.\ Impact of AI use on the project}
AI accelerated implementation, freeing time to focus on understanding concepts and interpreting results, and the predict-and-feedback dialogue helped internalize the concepts. At the same time, AI output is not always reliable, so all code, results, and claims were verified directly. Final responsibility rests with the team.

\end{document}
